% ═══════════════════════════════════════════════════════════════
% 第一章：引言
% ═══════════════════════════════════════════════════════════════

\section{引 言}

监管问询是中国资本市场信息披露制度中事中事后监管的核心工具。自2014年上海证券交易所和深圳证券交易所全面推行问询函制度以来，问询函的发放数量和覆盖范围持续扩大。2023年全年，两大交易所发出各类问询函超过1500份，覆盖超过三分之一的A股上市公司。对于上市公司而言，收到问询函不仅意味着信息披露瑕疵被公开曝光，还可能引发股价下跌\citep{chen2019supervision}、分析师评级下调\citep{wang2019analyst}以及投资者诉讼等一系列负面后果。因此，在问询函正式发出之前提前预判哪些公司可能"被问询"，对于投资风险管理和监管资源配置均具有重要的实践意义。

从学术研究的角度看，监管问询的可预测性问题蕴涵着重要的经济学含义。监管者的问询行为并非随机——它反映了监管者在信息不对称环境下对上市公司信息披露质量的系统性判断。已有文献从财务指标异常\citep{cassell2013review}、公司治理缺陷\citep{chen2019effective}和信息披露质量\citep{li2019inquiry}等角度考察了问询概率的驱动因素，但受限于数据和方法，现有研究存在三个方面的明显不足。

第一个不足是数据维度单一。现有预测研究几乎完全依赖结构化财务指标（如Beneish M-Score、资产负债率等），忽略了上市公司公告文本中蕴含的丰富语义信息。事实上，公告文本中的管理层语调、信息披露充分性、MD\&A前瞻性陈述等"软信息"维度，可能包含财务指标无法反映的重要风险信号。正如Liberti和Petersen（2019）在《经济文献杂志》的综述中所强调的，软信息和硬信息在风险评估中具有互补性，二者的融合能够带来比单一信息维度更优的预测表现\citep{liberti2019information}。

第二个不足是工具手段有限。大语言模型（Large Language Model, LLM）在语义理解方面的突破性进展为从非结构化文本中提取结构化信息提供了全新的技术手段\citep{huang2023finbert}。然而，LLM在监管问询预测中的应用尚属空白——现有文献既未系统检验LLM提取的语义特征对问询概率的增量预测能力，也未比较其相对于传统财务指标的相对重要性。

第三个不足是可解释性缺失。金融监管场景对AI模型的可解释性和可审计性有极高要求——投研和风控人员不仅需要知道"概率是多少"，更需要理解"为什么是这个概率"。现有预测模型多为"黑箱"式输出，缺乏从概率分数到风险特征再到原文证据的完整归因链。SHAP（Shapley Additive Explanations）等可解释性方法在金融领域的计量经济学应用中具有独特价值\citep{lu2023ml}，但其与大型语言模型生成式归因报告的融合在监管问询场景中尚缺乏系统探索。

本文在上述三个维度上做出贡献。在数据维度，本文构建了涵盖六大类约235个特征的多元信息融合框架（财务特征、LLM语义特征、市场特征、历史监管特征、知识图谱特征、时序衍生特征），将结构化与非结构化信息纳入统一分析框架。在方法维度，本文首次将大语言模型提取的文本语义风险特征引入监管问询概率预测，并采用Logistic回归、嵌套模型比较（Nested Model Comparison）\citep{delong1988comparing}和异构集成学习相结合的策略，系统检验软信息的增量预测价值及其在信息不对称不同公司间的异质性。在可解释性维度，本文构建了"SHAP数值归因—LLM生成式归因报告—原文证据追溯"三级可解释机制，并通过LLM-as-Judge框架对归因质量进行多维度量化评估\citep{zheng2024judge}。

基于中国A股上市公司约10万条公司-季度观测数据的实证分析表明：第一，传统财务异常指标对问询概率具有显著的预测能力（Logistic AUC 0.71）；第二，LLM提取的文本语义特征提供了显著且经济上重要的增量信息（全模型AUC 0.78，Delong检验p$<$0.01）；第三，软信息的增量价值在信息不对称更严重的小市值公司和低分析师覆盖公司中更加凸显；第四，LightGBM异质集成全模型的AUC达到0.79，Top-10\%高风险公司的覆盖率达38.6\%；第五，结合SHAP和LLM的三级可解释报告在LLM-as-Judge评估中获得86.3分的综合质量评价。

本文的研究发现对于完善事中事后监管制度、提升监管效率具有明确的政策含义：监管机构可参考本文融合多源信息的预测框架，将文本语义分析纳入日常监管工具箱，构建"财务指标+文本语义+市场信号"三位一体的智能化问询筛选系统。同时，可解释的归因报告机制能够帮助监管者和投资者更有效地利用AI预测结果，避免"黑箱"模型在金融监管场景中的潜在应用风险。

本文其余部分的结构安排如下：第二部分为制度背景与文献综述；第三部分为理论分析与研究假说；第四部分为数据、变量与描述性统计；第五部分为研究设计；第六部分为实证结果分析；第七部分为稳健性检验；第八部分为可解释性分析；第九部分为结论与政策建议。